% Generated draft: requires mathematical and source review.
\documentclass[11pt]{article}
\usepackage[margin=25mm]{geometry}
\usepackage{fontspec}
\setmainfont{Latin Modern Roman}
\usepackage{amsmath,amssymb,mathtools,unicode-math}
\setmathfont{Latin Modern Math}
\usepackage{enumitem,needspace,xurl,hyperref}
\setlength{\parindent}{0pt}
\setlength{\parskip}{5pt}
\setlength{\emergencystretch}{3em}
\newcommand{\N}{\mathbb{N}}
\newcommand{\Z}{\mathbb{Z}}
\newcommand{\Q}{\mathbb{Q}}
\newcommand{\R}{\mathbb{R}}
\newcommand{\C}{\mathbb{C}}
\newcommand{\F}{\mathbb{F}}
\newcommand{\A}{\mathbb{A}}
\newcommand{\PP}{\mathbb{P}}
\newcommand{\E}{\mathbb{E}}
\newcommand{\eps}{\varepsilon}
\newcommand{\dd}{\,\mathrm{d}}
\begin{document}
\section*{UnsolvedMath ID: 30003029}
\textbf{Source problem code:} catalogue entry
\subsection{Definitions and mathematical statement}
The supplied record specifies the Gaussian pair potential
\[
V_{\alpha}(r)=e^{-\alpha r^{2}},\qquad r\geq 0,
\]
and particle density \(1\) in Euclidean space. The dimension \(d\) of the ambient space \(\R^{d}\) is not specified. The question asks for the point configurations that minimize average energy, with particular emphasis on
\[
\alpha\geq \frac{4\pi}{e}.
\]
The intended parameter range outside this emphasized regime is not stated.

To display the intended minimization question without supplying missing hypotheses, let \(\mathcal C_{d,1}\) stand for the intended admissible class of density-one point configurations, and let \(\mathcal E_{\alpha}(X)\) stand for the intended average pair energy of \(X\) associated with \(V_{\alpha}\). The requested ground-state configurations would then form
\[
\mathcal G_{\alpha,d}
=
\left\{X\in\mathcal C_{d,1}:
\mathcal E_{\alpha}(X)\leq\mathcal E_{\alpha}(Y)
\text{ for every }Y\in\mathcal C_{d,1}\right\}.
\]
Here \(\mathcal C_{d,1}\) and \(\mathcal E_{\alpha}\) are placeholders for missing specifications, not fully defined objects supplied by the source record. In particular, the record does not define the density limit, the admissible configurations, the averaging procedure for the infinite system, the treatment of self-interactions, or boundary conditions. It does not specify whether minimization is over all point configurations, periodic configurations, or lattices. Consequently, the displayed formulation remains conditional on these choices.

The supplied dated literature assessment describes the system as unspecified, whereas the supplied statement explicitly names the Gaussian core model and its pair potential. This discrepancy remains unresolved; neither the missing specifications nor a literature conclusion are inferred from it.
\subsection{Short English statement}
Identify the density-one ground-state point configurations for the Gaussian core model with potential \(e^{-\alpha r^{2}}\), especially when \(\alpha\geq 4\pi/e\). The dimension, admissible configurations, and definition of average energy remain unspecified, so a complete mathematical formulation requires further source clarification.
\subsection{Sources}
\begin{itemize}
\item Local UnsolvedMath registry: supplied record with permanent ID \(30003029\), exact source problem code \texttt{catalogue entry}, title \emph{Ground States of the Gaussian Core Model}, in the Oberwolfach Reports Open Problems collection. The record supplies proposed year 2016 and names no individual proposer.
\item Supplied report citation: \emph{Lattices and Applications in Number Theory} (2016), Oberwolfach Reports, pp.\ 87--154, DOI \texttt{10.4171/owr/2016/3}. Source URL: \url{https://doi.org/10.4171/owr/2016/3}. The supplied background describes this report as documenting a mathematical workshop and its open-problem discussions; no problem-specific page or excerpt is supplied.
\item Supplied literature source: \emph{Oberwolfach Report 3/2016}, \url{https://doi.org/10.4171/owr/2016/3}. The exact supplied literature-check date is \texttt{2026-08-21}. Its assessment and open-status label are supplied dated curation, not independently verified here. The assessment says that the system, Hamiltonian, and admissible class are unspecified and that the exact claim cannot be classified more strongly; the supplied background requests independent verification of time-sensitive status and cited claims.
\item The registry labels the statement \texttt{corrected\_verified} and says that only safe notation, typography, and TeX-delimiter cleanup was applied, preserving all hypotheses, quantifiers, and subquestions. These are supplied registry qualifications, not an independent mathematical or source verification.
\end{itemize}
\end{document}
